\section{Introduction}
\label{sec:introduction}
Industrial maintenance and inspection are essential for ensuring the operational continuity and safety of manufacturing facilities. Traditionally, these tasks have relied on human expertise, which is subject to fatigue and inconsistency. With the advent of Industry 4.0, automated inspection systems leveraging deep learning have emerged as a robust alternative. 

Despite the progress, deploying automated systems in real-world industrial environments remains challenging. Standard computer vision models often struggle with (1) extreme perspective distortions and reflections on gauge surfaces, (2) the dense arrangement of identical indicators leading to matching drift, and (3) the diverse nature of inspection targets that range from simple LEDs to complex semantic signs. There is an urgent need for a system that provides high-precision diagnostic capability while remaining computationally efficient for robotic patrolling.

This paper introduces an integrated diagnostic system that combines state-of-the-art object detection with specialized geometric analysis modules and Vision-Language Models (VLM). The primary contributions of this work are as follows:
\begin{itemize}
    \item \textbf{Layered Modular Architecture}: We propose a four-layer framework (Orchestration, Perception, Diagnostics, Resource Management) that allows for independent scaling and optimization of each functional unit.
    \item \textbf{Hybrid Geometric-AI Diagnostics}: We integrate probabilistic deep learning with deterministic geometric refinement (Homography and Hough transforms) to achieve sub-2\% error rates in gauge reading under extreme viewing angles.
    \item \textbf{Prioritized Inference Strategy}: We implement a task-specific orchestration logic that selectively invokes intensive VLM reasoning only for complex "ETC" items, utilizing parallel execution to maintain real-time performance (120ms avg. latency).
    \item \textbf{Robust Empirical Validation}: We demonstrate the system's effectiveness through a comprehensive dataset from four distinct industrial zones, providing detailed qualitative analysis of 12 distinct edge-case scenarios.
\end{itemize}
